\documentclass[11pt]{article}

\usepackage[margin=1in]{geometry}
\usepackage{enumitem}
\usepackage{amsmath, amssymb}

\setlist[enumerate]{itemsep=0.6em, topsep=0.6em}

\begin{document}

\begin{center}
  {\Large \textbf{MED 520 Examination}}\\[0.25em]
  {\large \textbf{General Pathology and Disease Mechanisms -- Answer Key}}\\[0.5em]
\end{center}

\begin{enumerate}[label=\textbf{\arabic*.}]

  \item \textbf{(B)} Hypertrophy.\\
  Hypertrophy is an increase in cell size, often in response to increased workload. Example: cardiac 
  muscle hypertrophy in hypertension. In contrast, hyperplasia is an increase in cell number.

  \item \textbf{(C)} Caseous necrosis.\\
  Caseous necrosis has a distinctive cheese-like appearance and is characteristic of granulomatous 
  inflammation, particularly in tuberculosis. It combines features of coagulative and liquefactive 
  necrosis.

  \item \textbf{(C)} Neutrophils.\\
  Neutrophils (polymorphonuclear leukocytes) are the first responders in acute inflammation, arriving 
  within minutes to hours. They phagocytose bacteria and debris. Macrophages arrive later in acute 
  inflammation and dominate chronic inflammation.

  \item \textbf{(C)} AFP (Alpha-Fetoprotein).\\
  AFP is elevated in hepatocellular carcinoma and certain germ cell tumors. PSA is specific for 
  prostate cancer, CA-125 for ovarian cancer, and CEA is associated with colorectal and other cancers.

  \item \emph{Sample answer:}\\
  \textbf{Benign tumors}:
  \begin{itemize}
    \item \textbf{Growth}: Slow, expansile, often encapsulated
    \item \textbf{Invasion}: Do not invade surrounding tissues
    \item \textbf{Metastasis}: Do not metastasize
    \item \textbf{Differentiation}: Well-differentiated, resemble tissue of origin
    \item Example: Lipoma (benign adipose tumor)
  \end{itemize}

  \textbf{Malignant tumors}:
  \begin{itemize}
    \item \textbf{Growth}: Rapid, infiltrative
    \item \textbf{Invasion}: Invade adjacent tissues and destroy normal architecture
    \item \textbf{Metastasis}: Spread to distant sites via lymphatics and blood
    \item \textbf{Differentiation}: Poorly differentiated (anaplastic) to well-differentiated
    \item Example: Liposarcoma (malignant adipose tumor)
  \end{itemize}

  \textbf{Why metastasis is critical}: Metastasis is the primary cause of cancer-related mortality. 
  Local tumors can often be surgically removed, but metastatic disease is difficult to treat and 
  leads to multi-organ failure. The ability to metastasize definitively classifies a tumor as malignant.

  \item \emph{Sample answer:}\\
  \textbf{Pathophysiology}: Atherosclerosis is characterized by the formation of atherosclerotic 
  plaques (atheromas) in arterial walls, leading to stenosis, thrombosis, and ischemia.

  \textbf{Major risk factors}:
  \begin{itemize}
    \item Hyperlipidemia (especially LDL cholesterol)
    \item Hypertension
    \item Smoking
    \item Diabetes mellitus
    \item Age, male sex, family history
  \end{itemize}

  \textbf{Mechanisms}:
  \begin{enumerate}
    \item \textbf{Endothelial dysfunction}: Risk factors damage endothelium, increasing permeability 
    and promoting inflammation
    \item \textbf{Lipid accumulation}: LDL cholesterol infiltrates the intima and becomes oxidized (ox-LDL)
    \item \textbf{Inflammation}: Monocytes migrate into the intima, differentiate into macrophages, 
    and ingest ox-LDL, becoming foam cells
    \item \textbf{Smooth muscle proliferation}: Growth factors stimulate smooth muscle cell migration 
    and proliferation, forming a fibrous cap
    \item \textbf{Plaque formation}: Lipid core (foam cells, extracellular lipid) covered by fibrous 
    cap. Plaques can rupture, causing thrombosis
  \end{enumerate}

  Complications include myocardial infarction, stroke, and peripheral arterial disease.

  \item \emph{Sample answer:}\\
  \textbf{Apoptosis}: Programmed cell death characterized by controlled, energy-dependent dismantling 
  of the cell without inflammation.

  \textbf{Differences from necrosis}:
  \begin{itemize}
    \item \textbf{Apoptosis}: Cell shrinkage, chromatin condensation, membrane blebbing, formation 
    of apoptotic bodies, phagocytosis without inflammation
    \item \textbf{Necrosis}: Cell swelling, membrane rupture, release of contents, inflammatory response
  \end{itemize}

  \textbf{Two major pathways}:
  \begin{enumerate}
    \item \textbf{Intrinsic (mitochondrial) pathway}: Triggered by DNA damage, oxidative stress, 
    or growth factor withdrawal. Mitochondria release cytochrome c, activating caspases
    \item \textbf{Extrinsic (death receptor) pathway}: Triggered by binding of ligands (e.g., Fas 
    ligand, TNF) to death receptors, directly activating caspases
  \end{enumerate}

  \textbf{Role of caspases}: Caspases are proteases that execute apoptosis by cleaving cellular proteins, 
  leading to DNA fragmentation, cytoskeletal breakdown, and cell dismantling.

  \textbf{Significance}:
  \begin{itemize}
    \item \textbf{Development}: Eliminates unwanted cells (e.g., webbing between fingers)
    \item \textbf{Homeostasis}: Balances cell proliferation
    \item \textbf{Disease}: Excessive apoptosis causes atrophy (e.g., neurodegenerative diseases); 
    deficient apoptosis contributes to cancer
  \end{itemize}

  \item \emph{Sample answer:}\\
  \textbf{Carcinogenesis}: Multi-step process where normal cells progressively acquire genetic and 
  epigenetic alterations, leading to malignancy.

  \textbf{Multi-step model}:
  \begin{enumerate}
    \item \textbf{Initiation}: DNA damage by carcinogens (chemicals, radiation, viruses) creates mutations
    \item \textbf{Promotion}: Clonal expansion of initiated cells; mutations accumulate
    \item \textbf{Progression}: Additional mutations lead to invasion and metastasis
  \end{enumerate}

  \textbf{Key genetic changes}:
  \begin{itemize}
    \item \textbf{Oncogenes}: Mutated proto-oncogenes (e.g., RAS, MYC) promote uncontrolled proliferation. 
    Gain-of-function mutations; dominant (one mutated allele sufficient)
    \item \textbf{Tumor suppressor genes}: (e.g., TP53, RB) normally inhibit growth or promote apoptosis. 
    Loss-of-function mutations; recessive (both alleles must be inactivated, ``two-hit hypothesis'')
    \item \textbf{Apoptosis regulators}: (e.g., BCL-2 family) defects prevent elimination of damaged cells
    \item \textbf{DNA repair genes}: (e.g., BRCA1, BRCA2, mismatch repair genes) mutations increase 
    mutation rate, accelerating carcinogenesis
  \end{itemize}

  \textbf{Example}: Colorectal cancer progression: normal epithelium $\to$ adenoma (APC mutation) $\to$ 
  carcinoma (KRAS, TP53 mutations) $\to$ metastasis.

  Cancer results from the accumulation of multiple mutations affecting these pathways, conferring 
  hallmarks like sustained proliferation, evasion of apoptosis, angiogenesis, and metastatic capability.

\end{enumerate}

\end{document}
